
\setcounter{footnote}{-1}
\SGARefTarget{sga3:viii:chapter:viii}\chapter*{Expos\'e VIII\\Diagonalizable Groups}
\addcontentsline{toc}{chapter}{Expos\'e VIII. Diagonalizable Groups}
\markboth{EXPOS\'E VIII. DIAGONALIZABLE GROUPS}%
  {EXPOS\'E VIII. DIAGONALIZABLE GROUPS}

\begin{center}
  \textit{by A. Grothendieck}\footnote{Editor's note: version \SGARefLink{sga3:viii:definition:1:1}{1.1} of
  November 8, 2009. Additions were made in \SGARefLink{sga3:viii:theorem:1:2}{1.2}, \SGARefLink{sga3:viii:corollary:1:4}{1.4}, \SGARefLink{sga3:viii:paragraph:1-7}{1.7}, \SGARefLink{sga3:viii:theorem:3:1}{3.1}, \SGARefLink{sga3:viii:corollary:3:4}{3.4}, \SGARefLink{sga3:viii:remark:4:5:1}{4.5.1},
  \SGARefLink{sga3:viii:theorem:6:4}{6.4}, and \SGARefLink{sga3:viii:remark:6:8}{6.8}; 1.5.1 and \SGARefLink{sga3:viii:section:7}{Section~7} remain to be revised.}
\end{center}

\SGARefTarget{sga3:viii:section:1}\section*{1. Biduality}
\addcontentsline{toc}{section}{1. Biduality}

Let \(\mathscr C\) be a category, which as usual we identify with a full
subcategory of
\[
  \widehat{\mathscr C}
  =\operatorname{Hom}(\mathscr C^{\circ},\mathbf{Set})
\]
(cf. Expos\'e~I). Let \(I\) be a commutative group functor, that is, an
object of \(\widehat{\mathscr C}\) endowed with a commutative group structure
(cf. I, \SGARefLink{sga3:i:subsection:2-1}{2.1}).\footnote{Editor's note: the original has been expanded in what
follows.} For every \(X\in\operatorname{Ob}(\widehat{\mathscr C})\), the
object \(\operatorname{Hom}(X,I)\) carries the commutative group structure
induced by that of \(I\). For every group \(G\) in
\(\widehat{\mathscr C}\), put
\[
  D(G)=\operatorname{Hom}_{\mathrm{gr}}(G,I),
\]
the subobject of \(\operatorname{Hom}(G,I)\) defined, for every
\(S\in\operatorname{Ob}(\mathscr C)\), by
\SGARefTarget{sga3:viii:eq:1:x}
\begin{equation}
  \tag{x}
  D(G)(S)=\operatorname{Hom}_{S\text{-}\mathrm{gr}}(G_S,I_S),
\end{equation}
where \(G_S=G\times S\) and \(I_S=I\times S\) are regarded as
\(S\)-groups, that is, as groups in \(\widehat{\mathscr C}_{/S}\). Thus
\(D(G)\) is a \(\widehat{\mathscr C}\)-subgroup of
\(\operatorname{Hom}(G,I)\). In this way we obtain a contravariant functor
\(D\) from the category of \(\widehat{\mathscr C}\)-groups to the category
of commutative \(\widehat{\mathscr C}\)-groups.

The right-hand side of \(\SGARefLink{sga3:viii:eq:1:x}{(x)}\) may also
be interpreted as the subset of \(\operatorname{Hom}(G\times S,I)\)
consisting of morphisms \(G\times S\to I\) that are ``multiplicative in the
first argument \(G\).'' Moreover, the preceding formulas remain valid when
\(S\) is an arbitrary object of \(\widehat{\mathscr C}\), not necessarily
one arising from \(\mathscr C\).

Now take for \(S\) a group in \(\widehat{\mathscr C}\), denoted \(G'\). In
the left-hand side \(\operatorname{Hom}(G',D(G))\) of
\(\SGARefLink{sga3:viii:eq:1:x}{(x)}\), consider the subset
\(\operatorname{Hom}_{\mathrm{gr}}(G',D(G))\) consisting of morphisms that
respect the group structures of \(G'\) and \(D(G)\).

It corresponds to the subset of \(\operatorname{Hom}(G\times G',I)\)
consisting of morphisms that are multiplicative in each of the two
arguments---which one may call \emph{bilinear morphisms from
\(G\times G'\) to \(I\)}, or \emph{pairings of \(G\) and \(G'\) with
values in \(I\)}. We therefore have
\SGARefTarget{sga3:viii:eq:1:xx}
\begin{equation}
  \tag{xx}
  \operatorname{Hom}_{\mathrm{gr}}(G',D(G))
  \xrightarrow{\ \sim\ }
  \operatorname{Hom}_{\mathrm{bil}}(G\times G',I),
\end{equation}
functorially in the pair \((G,G')\). Since the right-hand side is symmetric
in \(G\) and \(G'\), this gives a functorial bijection
\SGARefTarget{sga3:viii:eq:1:xxx}
\begin{equation}
  \tag{xxx}
  \operatorname{Hom}_{\mathrm{gr}}(G',D(G))
  \xrightarrow{\ \sim\ }
  \operatorname{Hom}_{\mathrm{gr}}(G,D(G')).
\end{equation}
In other words, giving a group homomorphism \(G'\to D(G)\) ``amounts to the
same thing'' as giving a group homomorphism \(G\to D(G')\): both amount to
giving a pairing \(G\times G'\to I\).

Apply this with \(G'=D(G)\) and with the identity homomorphism
\(G'\to D(G)\). We obtain a canonical homomorphism
\SGARefTarget{sga3:viii:eq:1:xxxx}
\begin{equation}
  \tag{xxxx}
  G\longrightarrow D(D(G)).
\end{equation}

\medskip
\noindent\SGARefTarget{sga3:viii:definition:1:0}\textbf{Definition 1.0.}
\footnote{Editor's note: the numbering \SGARefLink{sga3:viii:definition:1:0}{1.0}, \SGARefLink{sga3:viii:proposition:1:0:1}{1.0.1}, and so on has been added
to make the definitions and statements contained here explicit.}
\textit{The group \(G\) is called \emph{reflexive} (relative to \(I\)) if
\SGARefLink{sga3:viii:eq:1:xxxx}{the preceding homomorphism} is an isomorphism. This implies, in particular,
that \(G\) is commutative.}

We therefore have:

\medskip
\noindent\SGARefTarget{sga3:viii:proposition:1:0:1}\textbf{Proposition 1.0.1.}
\textit{The functor \(D\) induces an anti-equivalence from the category of
reflexive \(\widehat{\mathscr C}\)-groups to itself.}

In particular, if \(G\) and \(H\) are reflexive groups, \(D\) induces an
isomorphism
\[
  \operatorname{Hom}_{\mathrm{gr}}(G,H)
  \xrightarrow{\ \sim\ }
  \operatorname{Hom}_{\mathrm{gr}}(D(H),D(G)).
\]
It is enough, in fact, that \(H\) be reflexive, as follows from
\(\SGARefLink{sga3:viii:eq:1:xxx}{(xxx)}\).

\medskip
\noindent\SGARefTarget{sga3:viii:definition:1:0:2}\textbf{Definition 1.0.2.}
\textit{As usual, a \(\mathscr C\)-group \(G\) is called \emph{reflexive} if
it is reflexive as a \(\widehat{\mathscr C}\)-group, regardless of whether
\(D(G)\) is representable.}

Thus \(D\) gives an anti-equivalence from the category of reflexive
\(\mathscr C\)-groups \(G\) for which \(D(G)\) is representable to itself.

\medskip
\noindent\SGARefTarget{sga3:viii:remark:1:0:3}\textbf{Remark 1.0.3.}
To conclude these generalities, note that formation of the dual \(D(G)\)
commutes with base extension, which therefore carries reflexive groups to
reflexive groups.

We shall henceforth work in the following setting. Let
\(\mathscr C=(\mathbf{Sch})_{/S}\) be the category of preschemes over \(S\),
and let \(I=\mathbf G_{m,S}\) be the ``multiplicative group over \(S\)''
(cf. Expos\'e~I). For every ordinary group \(M\), consider the constant
\(S\)-group \(M_S\). For every group prescheme \(J\) over \(S\), there is a
canonical isomorphism, functorial in \(M\) and \(J\) and compatible with base
extension,
\[
  \operatorname{Hom}_{S\text{-}\mathrm{gr}}(M_S,J)
  =\operatorname{Hom}_{\mathrm{gr}}(M,J(S)).
\]

Applying this with \(J=I=\mathbf G_{m,S}\) and with a variable \(S'\) over
\(S\), we obtain a functorial isomorphism
\SGARefTarget{sga3:viii:eq:1:0:4}
\begin{equation}
  \tag{1.0.4}
  D(M_S)(S')
  \xrightarrow{\ \sim\ }
  \operatorname{Hom}_{\mathrm{gr}}(M,\mathbf G_m(S')).
\end{equation}
We thus recover the functor already considered in I, \SGARefLink{sga3:i:subsection:4-4}{4.4}, also denoted
\(D_S(M)\), which is representable when \(M\) is commutative, since
\[
  D_S(M)=D(M_S)=\operatorname{Spec}\mathscr O_S(M),
\]
where \(\mathscr O_S(M)\) denotes the group algebra of \(M\) over
\(\mathscr O_S\). In general, moreover, \(D(M)\) is unchanged when \(M\)
is replaced by its abelianization, so no information is lost by assuming
that \(M\) is commutative.

\medskip
\noindent\SGARefTarget{sga3:viii:definition:1:1}\textbf{Definition 1.1.}
\textit{A group prescheme \(G\) over \(S\) is called \emph{diagonalizable}
if it is isomorphic to a scheme of the form}
\[
  D_S(M)=D(M_S)
  =\operatorname{Hom}_{S\text{-}\mathrm{gr}}(M_S,\mathbf G_m)
\]
\textit{for a suitable commutative group \(M\). It is called
\emph{locally diagonalizable} if every point of \(S\) has an open
neighborhood \(U\) such that \(G|_U\) is diagonalizable.}

\medskip
\noindent\SGARefTarget{sga3:viii:theorem:1:2}\textbf{Theorem 1.2.}
\textit{Let \(\Gamma\) be a constant commutative group scheme over \(S\);
that is, suppose that it is isomorphic to a group scheme of the form \(M_S\),
where \(M\) is an ordinary commutative group. Then \(\Gamma\) is reflexive:
the canonical homomorphism}
\[
  \Gamma\longrightarrow D(D(\Gamma))
\]
\textit{is an isomorphism.\footnote{Editor's note: the sentence that follows
has been added.} Hence the diagonalizable group \(D(M_S)\) is also
reflexive.}

In view of the definitions, this follows from \SGARefLink{sga3:viii:corollary:1:3}{the next statement}, applied to
a prescheme \(S'\) over \(S\).

\medskip
\noindent\SGARefTarget{sga3:viii:corollary:1:3}\textbf{Corollary 1.3.}
\textit{Let \(G=D(M_S)\). Every homomorphism of \(S\)-groups}
\[
  u:G\longrightarrow\mathbf G_{m,S}
\]
\textit{is defined by a unique section of \(M_S\) over \(S\), equivalently
by a unique locally constant map from \(S\) to \(M\).}

\emph{Proof.}
By definition,
\[
  \mathbf G_{m,S}
  =\operatorname{GL}(1)_S
  =\operatorname{Aut}_{\mathscr O_S\text{-}\mathrm{mod}}(\mathscr O_S).
\]
Thus giving a group homomorphism \(G\to\mathbf G_{m,S}\) is equivalent to
giving \(\mathscr O_S\) the structure of a
\(G\)-\(\mathscr O_S\)-module compatible with its natural
\(\mathscr O_S\)-module structure (cf. I, \SGARefLink{sga3:i:subsection:4-7}{4.7}). By I, \SGARefLink{sga3:i:proposition:4-7-3}{4.7.3}, this is also
equivalent to giving an \(M\)-grading on \(\mathscr O_S\), that is, a direct
sum decomposition
\[
  \mathscr O_S=\bigoplus_{m\in M}\mathscr L_m.
\]
It is well known that a direct summand of a locally free module of finite
type is locally free of finite type. Hence, near each point of \(S\), every
\(\mathscr L_m\) is either zero or free of rank one; in the latter case it
identifies with \(\mathscr O_S\) on that neighborhood. Let \(S_m\) be the
open subset of \(S\) on which the latter alternative occurs. The direct-sum
decomposition shows that the \(S_m\) cover \(S\) and are pairwise disjoint.

Consequently, giving a group homomorphism \(G\to\mathbf G_{m,S}\) is
equivalent to giving a decomposition of \(S\) as a union of pairwise
disjoint open subsets \(S_m\), \(m\in M\), hence to giving a locally
constant map from \(S\) to \(M\). This proves \SGARefLink{sga3:viii:corollary:1:3}{Corollary~1.3} and therefore
\SGARefLink{sga3:viii:theorem:1:2}{Theorem~1.2}. \qed
